\chapter{Écrire une IA}

Le fichier de ce chapitre :
\nkfile{Applications/ConquerorLab/exemples/ai/IAMinimale.cpp}. Il tient en une
idée : \textbf{demander les coups légaux au moteur, en prendre un au hasard.}
C'est nul, et c'est le but.

\section{Pourquoi commencer par une IA nulle}

Une IA aléatoire est le \textbf{plancher de mesure}. Sans elle, vous n'avez aucun
point de comparaison : une IA qui gagne 70 \% ne veut rien dire tant qu'on ne sait
pas contre quoi.

Le premier run du banc d'essai livre déjà un enseignement de cette nature :

\begin{quote}\small
Le glouton bat l'aléatoire 29--6. La position se juge, l'évaluation matérielle
mord.
\hfill\textsc{ConquerorLab/README.md}
\end{quote}

C'est une information réelle sur le jeu, obtenue avec deux IA triviales. Une IA
qui ne bat pas l'aléatoire ne vaut rien ; une qui la bat 55 \% du temps ne vaut
pas beaucoup plus.

\section{Le point de conception central}

\begin{nkcode}{include/Conqueror/ConquerorAIABI.h:14-22}
POINT DE CONCEPTION CENTRAL
L'IA ne connait PAS l'implementation du moteur de regles : elle recoit une
NkcRulesVTable et un handle opaque. Consequences voulues :

  1. Le sujet A2 demarre en semaine 1 sur un module bouchon, sans attendre A1.
  2. Quand A1 change son code interne, l'IA continue de fonctionner.
  3. Les deux ne peuvent pas diverger sur les regles : il n'en existe qu'une
     implementation, et l'IA la traverse.
\end{nkcode}

Concrètement, votre IA ne sait pas ce qu'est une duplication. Elle ne connaît ni
le plateau, ni la topologie, ni les paramètres. Elle demande :

\begin{nkcode}{exemples/ai/IAMinimale.cpp — V\_ChooseMove}
// TOUT passe par la table : on ne sait pas ce qu'on joue, on demande.
NkcMove      moves[kMoveCap];
const uint32 total = R->GenerateLegalMoves(ri, st, moves, kMoveCap);
const uint32 n     = total < kMoveCap ? total : kMoveCap;
if (n == 0) return 0;
\end{nkcode}

\begin{nkretenir}
Vous pouvez écrire une IA \textbf{avant} que les règles existent. Vous pouvez la
tester contre trois moteurs différents sans la recompiler. Et le jour où A1
change sa représentation interne, vous ne le saurez même pas.
\end{nkretenir}

\section{Les neuf fonctions}

\begin{center}
\begin{tabular}{@{}p{4.6cm}p{2.1cm}p{6.6cm}@{}}
\toprule
\textbf{Fonction} & \textbf{Obligatoire} & \textbf{Rôle} \\
\midrule
\texttt{Create} / \texttt{Destroy} & oui & cycle de vie \\
\texttt{Configure} & oui & palier, budget, graine — avant le 1\textsuperscript{er}
  \texttt{ChooseMove} \\
\texttt{ChooseMove} & oui & \textbf{le cœur} \\
\texttt{GetParamsSchemaJson} / \texttt{SetParam} & non & réglages ; renvoyer
  \texttt{""} est légal \\
\texttt{OnMovePlayed} & non & un coup a réellement été joué \\
\texttt{Reset} & non & nouvelle partie : vider la mémoire \\
\texttt{GetDebugJson} & non & ce que l'atelier pourra afficher \\
\bottomrule
\end{tabular}
\end{center}

\section{\texttt{ChooseMove}, ligne à ligne}

\begin{nkcode}{exemples/ai/IAMinimale.cpp — V\_ChooseMove, entree}
int32 V_ChooseMove(NkcAI self, const NkcRulesVTable *R, NkcRules ri,
                   const NkcState st, NkcAIResult *out) {
    AI *a = static_cast<AI *>(self);
    if (!a || !R || !ri || !st || !out) return 0;

    std::memset(out, 0, sizeof(*out));
    out->move.targetLevel = -1;
    out->move.powerId     = -1;
\end{nkcode}

\begin{nkpiege}[Le même piège que pour les règles]
\texttt{NkcMove} se compare \textbf{octet à octet}. Un champ laissé au hasard rend
votre coup illégal aux yeux du moteur.

Ce n'est pas strictement nécessaire quand on recopie un coup issu de
\texttt{GenerateLegalMoves} — il est déjà propre — mais une IA qui
\emph{construit} un coup au lieu de le choisir doit absolument le faire.
\end{nkpiege}

\begin{nkcode}{exemples/ai/IAMinimale.cpp — V\_ChooseMove, choix}
    const uint32 pick = static_cast<uint32>(NextRand(a->rng) % n);
    out->move         = moves[pick];
    out->simulations  = 1;
    out->depthReached = 0;
    out->scoreMilli   = 0;
    return 1;
}
\end{nkcode}

Le retour est binaire : \textbf{1 = j'ai un coup, 0 = je n'en ai aucun}. Dans le
second cas l'atelier joue PASSER lui-même.

\subsection{Le compte-rendu}

\texttt{NkcAIResult} n'est pas décoratif : le panneau \emph{Joueurs} l'affiche
après chaque réflexion.

\begin{nkcode}{include/Conqueror/ConquerorAIABI.h:78-86}
struct NkcAIResult {
        NkcMove move;
        int32   scoreMilli   = 0;  // -1000 = perdue, +1000 = gagnee
        uint32  simulations  = 0;  // simulations/noeuds reellement faits
        uint32  elapsedMs    = 0;
        uint32  depthReached = 0;
        uint8   hitBudget    = 0;  // 1 si coupee par le budget
        uint8   _pad[3]      = {};
};
\end{nkcode}

\begin{nkretenir}[\texttt{scoreMilli} est en MILLIÈMES]
Aucun flottant à la frontière, même pour un score. Même raison qu'ailleurs : deux
plateformes doivent produire le même nombre.
\end{nkretenir}

\section{Le PRNG : porté par l'instance}

\begin{nkcode}{exemples/ai/IAMinimale.cpp — NextRand}
// PRNG xorshift64* PORTE PAR L'INSTANCE, jamais global. Sans cela, deux
// instances d'IA jouant en parallele dans une campagne piochent dans le
// meme flux, et deux parties de meme graine cessent d'etre reproductibles.
uint64 NextRand(uint64 &s) {
    s ^= s >> 12;
    s ^= s << 25;
    s ^= s >> 27;
    return s * 2685821657736338717ull;
}
\end{nkcode}

La graine arrive par \texttt{Configure}, et l'atelier la dérive de
\emph{(partie, tour, joueur)} :

\begin{nkcode}{src/ConquerorLab/NkcSession.h — StartThinking}
// Graine derivee de (partie, tour, joueur) : deux parties de meme graine
// rejouent a l'identique, et deux tours ne partagent pas le meme flux.
c.seed = mSeed ^ (static_cast<uint64>(mView.turn) * 0x9E3779B97F4A7C15ull) ^
         (static_cast<uint64>(mView.current) + 1ull);
\end{nkcode}

\begin{nkpiege}[Un \texttt{static} dans votre IA, et la campagne ment]
Une campagne fait tourner \textbf{plusieurs instances de votre IA en parallèle},
une par worker. Un compteur \texttt{static}, un cache global, un générateur
partagé : et vos parties cessent d'être indépendantes.

Le résultat reste \emph{plausible} — c'est le pire des cas, parce que vous ne le
verrez pas. Déclarez-le franchement dans votre fabrique :
\texttt{info.isThreadSafe = 0} si vous n'en êtes pas sûr.
\end{nkpiege}

\section{Le thread, et le budget}

\begin{nkcode}{include/Conqueror/ConquerorAIABI.h:24-32}
CONTRAINTE DE THREAD
ChooseMove est appelee depuis un THREAD WORKER, jamais depuis le thread de
rendu. Elle doit :
  - ne toucher que l'etat qu'on lui passe (c'est un CLONE prive) ;
  - ne poser aucun verrou sur une ressource partagee ;
  - respecter `budgetMs` : au-dela, elle rend le meilleur coup trouve
    jusque-la. Depasser le budget fige l'interface -- c'est un echec.
\end{nkcode}

Côté atelier, l'isolement est total :

\begin{quote}\small
L'IA réfléchit sur son propre moteur. Plutôt que de partager l'instance de règles
avec le thread de rendu — où l'utilisateur peut à tout moment bouger un
paramètre — le thread possède SA PROPRE instance, synchronisée avant chaque
réflexion, et son propre état, transféré par
\texttt{SerializeState}/\texttt{DeserializeState}. Zéro verrou, zéro course.
\hfill\textsc{NkcSession.h}
\end{quote}

L'état que vous recevez est \textbf{à vous}. Vous pouvez le cloner autant que vous
voulez. Mais vous ne devez \textbf{pas} le muter directement — vous passez par
\texttt{rules->ApplyMove} sur vos clones.

\subsection{Respecter le budget}

\texttt{IAMinimale} répond instantanément, la question ne se pose pas. Dès que
vous faites une recherche, elle se pose à chaque nœud.

\begin{nkcode}{Idiome — ecrit pour ce cours, d'apres ConquerorAIABI.h:113}
// Au debut de ChooseMove : noter l'instant de depart (l'IA choisit sa
// source de temps ; le contrat n'en impose aucune).
// Dans la boucle de recherche, tous les N noeuds :
if (a->cfg.budgetMs != 0 && ElapsedMs(start) >= a->cfg.budgetMs) {
    out->hitBudget = 1;
    break;                       // on rend le MEILLEUR coup trouve jusque-la
}
\end{nkcode}

Deux règles :

\begin{itemize}
  \item \texttt{budgetMs == 0} signifie \textbf{illimité}. Ne le traitez pas
        comme « zéro milliseconde » ;
  \item \textbf{ne testez pas le temps à chaque nœud} : un appel horloge par nœud
        coûte plus cher que le nœud. Tous les 1024 nœuds suffit.
\end{itemize}

\begin{nkpiege}[Ce que « figer l'interface » veut dire concrètement]
Le panneau \emph{Joueurs} affiche \texttt{hitBudget} sous la forme « coupée par le
budget ». Mais surtout : quand vous cliquez \emph{Nouvelle partie} pendant une
réflexion, l'atelier \textbf{attend} que votre \texttt{ChooseMove} rende la main.
Une IA qui ignore son budget rend le bouton mou pendant tout ce temps.
\end{nkpiege}

\section{Les paliers de difficulté}

\begin{nkcode}{include/Conqueror/ConquerorAIABI.h:41-51}
/// Cinq paliers. Reperes, pas implementation imposee : c'est au module de
/// decider comment il module sa force (budget de simulations, profondeur,
/// bruit, erreurs volontaires).
enum class NkcDifficulty : uint8 {
    Easy = 0, Normal = 1, Hard = 2, Expert = 3, Apex = 4, Count = 5
};
\end{nkcode}

L'IA de référence les utilise comme trois stratégies distinctes : aléatoire pur
en \emph{Facile}, glouton en \emph{Normal}, minimax alpha-bêta de profondeur 2, 3
puis 4 pour les trois derniers.

Si votre IA ne distingue qu'un seul niveau, dites-le :
\texttt{info.difficultyCount = 1}.

\section{Évaluer une position sans connaître les règles}

C'est le problème intéressant du sujet A2. Voici comment l'IA de référence s'y
prend — et c'est exactement ce qu'il faut remettre en cause.

\begin{nkcode}{modules/ai/ConquerorAIRef.cpp — Evaluate}
int32 mine = 0, theirs = 0;
for (uint8 p = 0; p < v.playerCount; ++p) {
    if (p == me) mine += v.totemCount[p];
    else         theirs += v.totemCount[p];
}

// Mobilite du joueur au trait, bornee pour rester bon marche.
NkcMove      buf[64];
const uint32 n        = R->GenerateLegalMoves(ri, st, buf, 64);
int32        mobility = static_cast<int32>(n > 64 ? 64 : n);
if (v.current != me) mobility = -mobility;

return ai->wMaterial * (mine - theirs) + ai->wMobility * mobility;
\end{nkcode}

Deux termes seulement — le \textbf{matériel} et la \textbf{mobilité} — et deux
poids réglables. Tout ce qu'une IA peut savoir vient de \texttt{NkcStateView} et
de \texttt{GenerateLegalMoves}.

\begin{nkretenir}[Le vrai sujet de A2]
Cette évaluation est délibérément grossière. « Le glouton bat l'aléatoire 29--6 »
dit que le matériel mord ; il ne dit pas que c'est le bon critère.

Une position avec beaucoup de totems mais peu de mobilité est-elle bonne ? Un
totem isolé au centre vaut-il un totem en bord de plateau ? \textbf{Ce sont des
questions que seule la mesure tranche} — chapitre 6.
\end{nkretenir}

\section{\texttt{OnMovePlayed} : ce qui distingue un MCTS d'une recherche jetable}

Une recherche naïve jette tout son travail à chaque tour. Un MCTS conserve son
arbre : quand un coup est joué, il \emph{descend} dans le sous-arbre
correspondant et repart de là. \texttt{OnMovePlayed} est le signal qui le lui
permet, \texttt{Reset} celui qui lui dit de tout jeter.

\begin{nkpiege}[Attention en campagne]
Une campagne réutilise la même instance d'IA d'une partie à l'autre. Sans
\texttt{Reset} correct, l'arbre de la partie précédente contamine la suivante — et
vos mesures deviennent dépendantes de l'ordre des parties.

Le lanceur de campagne appelle \texttt{Configure} puis \texttt{Reset} au début de
chaque partie ; à vous d'honorer le second.
\end{nkpiege}

\section{Votre première IA, en quatre gestes}

\begin{enumerate}
  \item copiez \nkfile{exemples/ai/IAMinimale.cpp} vers
        \nkfile{travail/ai/mon\_ia.cpp} ;
  \item changez le nom dans \texttt{FillFactory} ;
  \item sauvegardez, attendez une seconde, ouvrez le panneau \emph{Joueurs} :
        votre IA est dans la liste des pilotes ;
  \item mettez-la sur le siège 1, l'IA de référence sur le siège 0, et lancez une
        campagne de 200 parties.
\end{enumerate}

Vous venez de mesurer votre plancher. Tout ce que vous ferez ensuite se compare à
ce chiffre.

\begin{nkexo}[1 — Le glouton]
Transformez \texttt{IAMinimale} en glouton : pour chaque coup légal, clonez
l'état, appliquez le coup, comptez vos totems, gardez le meilleur. Vous aurez
besoin de \texttt{R->CreateState}, \texttt{CloneState}, \texttt{ApplyMove},
\texttt{GetView}, \texttt{DestroyState}. Mesurez-le contre l'aléatoire sur 200
parties.
\end{nkexo}

\begin{nkexo}[2 — Le poids réglable]
Exposez un paramètre \texttt{poids\_mobilite} via \texttt{GetParamsSchemaJson}
(même format que pour les règles). Vérifiez qu'il apparaît dans l'interface, puis
faites-le varier et mesurez : y a-t-il une valeur optimale, ou le résultat est-il
plat ?
\end{nkexo}

\begin{nkexo}[3 — Le budget]
Ajoutez une recherche à profondeur 3, puis un test de budget tous les 1024 nœuds.
Réglez le budget à 5 ms dans le panneau \emph{Joueurs} et vérifiez que
\texttt{hitBudget} passe bien à 1. Que se passe-t-il si vous oubliez le test ?
\end{nkexo}

\begin{nkexo}[4 — Le piège du \texttt{static}]
Ajoutez un \texttt{static int compteur = 0;} dans votre \texttt{ChooseMove} et
faites-le intervenir dans le choix du coup. Lancez une campagne sur 8 threads,
deux fois, avec la même graine. Les résultats sont-ils identiques ? Expliquez.
\end{nkexo}
